%% kapa: exam=1 round=30 year=2019 period=1 form=A subject=civil_law subject_ko=민법 seq=26 number=26
%% source: Q-Net 감정평가사 기출문제 (cst00309) artlSeq=5205205 file=extracted/1008/5205205/2019년 제30회 감정평가사 1차시험 문제지/2019년 제30회 감정평가사 1차시험 문제지/2019년도 제30회 감정평가사 1차 1교시 A형.hwp
%% conversion: kordoc 4.24.51 -> tools/gen1.py (auto); review=proofread-sonnet; flags=none
\begin{question}{26}
선의취득에 관한 설명으로 옳지 않은 것은? (다툼이 있으면 판례에 따름)
\begin{choices}
  \item 등록에 의하여 소유권이 공시되는 자동차는 동산이라 하더라도 선의취득의 대상이 되지 않는다.
  \item 수분양자로서의 지위를 내용으로 하는 연립주택의 입주권은 선의취득의 대상이 될 수 없다.
  \item 채무자 이외의 자의 소유에 속하는 동산의 경매절차에서 그 동산을 경락받아 경락대금을 납부하고 이를 인도받은 경락인은 특별한 사정이 없는 한 소유권을 선의취득할 수 있다.
  \item 선의취득이 인정되기 위해서는 양도인이 무권리자인 점을 제외하면 아무런 흠이 없는 거래행위이어야 한다.
  \item 현실인도뿐만 아니라 점유개정의 방법으로 양수인이 동산의 점유를 취득한 경우에도 선의취득이 인정된다.
\end{choices}
\end{question}
\answer{5}
